% ============================================================
%  第六章  辛普森悖论、混杂与可压缩性（第二部分：6.4—6.6）
% ============================================================

\section{稳定无偏与偶然无偏}

\subsection{动机}

上一例中关联准则的失效，要求我们重新审视 \textup{(6.10)} 所定义的混杂与无偏性概念。在例 6.3.3 中，$X$ 与 $Y$ 之所以被归为无混杂，是因为通过令 $r=-\alpha\gamma$，我们能使由 $r$ 代表的虚假关联抵消由 $Z$ 中介的那一份。在实践中，这种完美的抵消将是特定于某种研究条件之奇特组合的偶然事件；当问题参数（即 $\alpha$、$\gamma$ 与 $r$）发生轻微变化时——比如研究在不同地点或不同时间重复进行时——它便不会持续。相反，例 6.3.1 中的无混杂条件并不表现出这种易变性。在该例中，无论教育与暴露之间的联系强度如何，也无论教育与年龄如何影响病人拥有的车型，\textup{(6.10)} 所表达的无偏性都将继续成立。我们把这类无偏性称为\emph{稳定的}，因为它对参数变化稳健，只要模型中因果联系的构型不变，它就保持不变。

鉴于稳定无偏与偶然无偏的这一区别，我们需要重新审视：若一个准则把仅由偶然抵消而造成无混杂的情形误判为混杂，我们是否应认为它不充分？更根本地，鉴于这些情形中 \textup{(6.10)} 得以成立的条件颇为脆弱，我们是否应坚持把这类奇特情形纳入无偏性的定义？这些问题的一部分答案取决于选择，但有充分证据表明，我们关于混杂的直觉是由稳定无偏性的考量、而非仅偶然的考量所驱动的。否则我们如何解释：一代代流行病学家与生物统计学家会倡导那些在涉及偶然抵消的情形中失效的混杂准则？从实用角度看，未能识别偶然无偏的情形不应在观测研究中引入可察觉的误差，因为这些情形转瞬即逝，很可能在略有不同的条件下被后续研究否证。\footnote{正如我们在例 6.3.3 中所见，任何能够识别这类情形的统计检验，都要求测量 $T$ 中的所有变量。}

假定我们准备只把无偏性对参数变化保持稳健的情形归为无偏，那么仍有两个问题：(1) 如何给“稳定无偏性”这一新概念一个形式的、非参数的表述？(2) 是否有实用的统计准则来检验稳定无偏性？两个问题都可借助结构模型回答。

\subsection{形式定义}

第三章描述了一个称为“后门准则”的图形准则，用以在因果图中识别无偏性的条件。\footnote{关于后门准则在流行病学中之应用的浅近介绍，见 Greenland et al. (1999a)。}在不对（被测量协变量）作调整的简单情形中，该准则断言：若 $X$ 与 $Y$ 之间每条含有指向 $X$ 之箭头的路径，也含有对头对头指向的箭头（如图 6.3），则 $X$ 与 $Y$ 无混杂；只要图中缺失的连线代表相应变量之间不存在因果联系，该准则便有效。由于缺失连线中蕴含的因果假设如此明确，后门准则有两个显著特征。第一，不需要任何统计信息；图的拓扑结构便足以可靠地判定一个效应是否无混杂（在定义 6.2.1 的意义上），以及在混杂存在时对一组变量作调整是否足以消除混杂。第二，任何满足后门准则的模型，事实上会对无穷多个模型（或情形）类满足 \textup{(6.10)}，而该类中的每个模型都是通过给图中因果联系赋予不同参数而生成的。

为说明这一点，考虑图 6.3 所示的图。后门准则将把数对 $(X,Y)$ 识别为无混杂，因为唯一以指向 $X$ 的箭头结尾的路径是经过 $(X,E,Z,A,Y)$ 的那条，而这条路径在 $Z$ 处含有两个对头指向的箭头。此外，由于该准则仅基于图形关系，显然无论图中箭头所代表的因果关系的强度或类型如何，$(X,Y)$ 都将继续被归为无混杂。相反，考虑例 6.3.3 的图 6.4，其中有两条路径以指向 $X$ 的箭头结尾。由于这些路径都不含对头箭头，后门准则不会把 $X$ 对 $Y$ 的效应归为无混杂，从而承认：等式 $r=-\alpha\gamma$（若成立）并不代表稳定的无偏情形。

后门准则对因果假设的脆弱性，可在图 6.3 的语境中演示。假定研究者怀疑变量 $Z$（车型）对结局变量 $Y$ 有某种影响。这就相当于在图中添加一条从 $Z$ 到 $Y$ 的箭头，把该情形归为混杂，并提示对 $E$（或 $\{A,Z\}$）作调整。然而，若由于该研究特定的实验条件，$Z$ 事实上对 $Y$ 没有影响，则无须任何调整。确实，后门准则所建议的调整不会引入偏倚，但若它要求在无混杂情形中作多余的测量，这种调整可能代价高昂。\footnote{表面上，Stone–Robins 准则似乎能正确识别此情形中混杂的缺失，因为它基于实际生成数据的概率分布中成立的关联（按该分布，$\{E,Z\}$ 应在给定 $\{A,X\}$ 时与 $Y$ 独立）。然而，这些关联无助于判断是否可以省去某些测量；这类决定必须在采集数据之前作出，因而必须依赖关于条件关联之消失的主观假设。此类假设通常由因果知识、而非关联知识支持（见 \S 1.3）。}考虑到 (i) 手头的因果信息（即 $Z$ 可能影响 $Y$），以及 (ii) 我们对确保稳定无偏性的坚持——即在一切与手头信息相容的情形中避免偏倚——所增加的代价是正当的。

为形式地区分稳定无偏与偶然无偏，我们使用下述一般定义。

\begin{definition}[稳定无偏性]
设 $A$ 是关于数据生成过程的一组假设（或限制），$C_A$ 是满足 $A$ 的一类因果模型。若 $P(y\mid \dox{x})=P(y\mid x)$ 在 $C_A$ 中的每个模型 $M$ 中都成立，则称 $X$ 对 $Y$ 的效应估计在给定 $A$ 时是\emph{稳定无偏的}。相应地，我们说数对 $(X,Y)$ 在给定 $A$ 时是稳定无混杂的。
\end{definition}

用于指定因果模型的假设，既可以是参数的，也可以是拓扑的。例如，社会科学与经济学中使用的结构方程模型，通常受线性与正态性假设的限制。在这种情况下，$C_A$ 将由给方程中未指定的参数以及误差项协方差矩阵赋予不同数值所生成的全部模型构成。当我们仅指定因果图的拓扑结构，而让误差分布与方程的函数形式保持未定，就出现了更弱的非参数假设。我们现在探究这些非参数假设的统计意涵。

\begin{definition}[结构稳定无混杂]
设 $A_D$ 是嵌入因果图 $D$ 的假设集。若 $P(y\mid \dox{x})=P(y\mid x)$ 在 $D$ 的每个参数化中都成立，则称 $X$ 与 $Y$ 在给定 $A_D$ 时是\emph{稳定无混杂的}。所谓“参数化”，我们指给图的各条连线赋予函数、并给图中背景变量赋予先验概率。
\end{definition}

嵌入因果图的假设的明确解释见第三章与第五章。简言之，若 $D$ 是与因果模型相联系的图，则：
\begin{enumerate}[label=\arabic*.,leftmargin=1.8em]
\item 每条缺失的箭头（比如 $X$ 与 $Y$ 之间）代表如下假设：一旦我们作干预并保持 $Y$ 的父节点固定，$X$ 对 $Y$ 没有影响；
\item 每条 $X$ 与 $Y$ 之间缺失的双向连线代表如下假设：除 $D$ 中所示者外，$X$ 与 $Y$ 没有共同原因。
\end{enumerate}

只要图 $D$ 无环，后门准则就给出给定 $A_D$ 时稳定无混杂的一个充分必要检验。在不对协变量作调整的简单情形中，该准则归结为：$X$ 与 $Y$ 不存在（被观测或潜在的）共同祖先。\footnote{口语中的“共同祖先”应排除那些除经由 $X$ 之外与 $Y$ 没有其他联系的节点（例如图 6.3 中的节点 $E$），并包括代表相关误差的潜在节点。例如在图 6.4 中，应理解 $X$ 与 $Y$ 有两个共同祖先：第一个是 $Z$，第二个是造成 $X$ 与 $Y$ 之间双向弧的（隐含的）潜在变量（即 $\varepsilon_1$ 与 $\varepsilon_2$ 之间的相关）。}于是我们有下一个定理。

\begin{theorem}[共同原因原则]
设 $A_D$ 是嵌入无环因果图 $D$ 的假设集。$X$ 与 $Y$ 在给定 $A_D$ 时稳定无混杂，当且仅当 $X$ 与 $Y$ 在 $D$ 中没有共同祖先。
\end{theorem}

\noindent\textbf{定理 6.4.3 的证明}\quad “若”那一半由后门准则的有效性（定理 3.3.2）推出。“仅当”那一半要求构造一个特定模型，使得只要 $X$ 与 $Y$ 在 $D$ 中有共同祖先，\textup{(6.10)} 就被违反。用线性模型与 Wright 的路径系数规则容易做到。$\blacksquare$

定理 6.4.3 给出了稳定无混杂的一个充分必要条件，而不必援引统计数据，因为它完全依赖图中嵌入的信息。当然，图本身具有可检验的统计意涵（\S 1.2.3 与 \S 5.2.1），但这些检验并不能唯一地确定图（见第二章与 \S 5.2.3）。

然而，假定我们并不掌握构造因果图所需的全部信息，而只知道对每个变量 $Z$，是否可以安全地假定 $Z$ 对 $Y$ 没有影响，以及 $X$ 对 $Z$ 没有影响。现在的问题是：这些更为有限的信息连同统计数据，是否足以把数对 $(X,Y)$ 判为或判非稳定无混杂？答案是肯定的。

\subsection{稳定无混杂的操作检验}

\begin{theorem}[稳定无混杂的准则]
设 $A_Z$ 表示如下假设：(i) 数据由某个（未指定的）无环模型 $M$ 生成；(ii) $Z$ 是 $M$ 中一个不受 $X$ 影响、但可能影响 $Y$ 的变量。\footnote{所谓“可能影响 $Y$”，我们指：$A_Z$ 不包含“$Z$ 不影响 $Y$”这一假设。换言之，与 $M$ 相联系的图必须含有从 $Z$ 到 $Y$ 的有向路径。}若定义 6.2.2 的两个关联准则 (U\textsubscript{1}) 与 (U\textsubscript{2}) 都被违反，则 $(X,Y)$ 在给定 $A_Z$ 时不是稳定无混杂的。
\end{theorem}

\noindent\textbf{定理 6.4.4 的证明}\quad 只要 $X$ 与 $Y$ 稳定无混杂，定理 6.4.3 就排除了与底层模型相联系之图中 $X$ 与 $Y$ 存在共同祖先的可能。而共同祖先的不存在，又蕴含着：只要 $Z$ 满足 $A_Z$，则 (U\textsubscript{1}) 或 (U\textsubscript{2}) 之一必被满足。这是用 d-分离规则（\S 1.2.3）读取图所蕴含的条件独立关系的结果。\footnote{它也可由 Robins (1997) 的定理 7(a) 推出。}$\blacksquare$

定理 6.4.4 意味着，传统关联准则 (U\textsubscript{1}) 与 (U\textsubscript{2}) 可以用于一个关于稳定无混杂的简单操作检验；这一检验不要求我们知道领域中变量的因果结构，甚至不要求枚举相关变量的集合。只要找到任何一个满足 $A_Z$ 且违反 (U\textsubscript{1}) 与 (U\textsubscript{2}) 的变量 $Z$，就足以把 $(X,Y)$ 判为非稳定无混杂（尽管在研究所处的特定实验条件下，$(X,Y)$ 也许是偶然无混杂的）。

定理 6.4.4 传达了统计关联与混杂之间的一种形式联系，而这种联系并不基于封闭世界假设。\footnote{据我所知，文献中没有另一种这样的联系。}值得注意的是，这一联系可以在如此之弱的一组附加假设下形成：只需“某变量可能影响 $Y$ 且不受 $X$ 影响”这一定性假设，就足以产生一个关于稳定无混杂的必要统计检验。

\section{混杂、可压缩性与可交换性}

\subsection{混杂与可压缩性}

定理 6.4.4 还建立了混杂与“可压缩性”之间的形式联系——后者是这样一种准则：某种关联度量在省略某些变量时保持不变。

\begin{definition}[可压缩性]
设 $g[P(x,y)]$ 是任何度量联合分布 $P(x,y)$ 中 $Y$ 与 $X$ 之间关联的泛函。\footnote{泛函是把实数赋予给定函数集中任一函数的指派。例如，均值 $E(X)=\sum_x xP(x)$ 是一个泛函，因为它把实数 $E(X)$ 赋予每个概率函数 $P(x)$。}若
\[
E_z g[P(x,y\mid z)] = g[P(x,y)],
\]
则称 $g$ 关于变量 $Z$ \emph{可压缩}。
\end{definition}

不难证明，若 $g$ 代表 $P(y\mid x)$ 的任何线性泛函——例如风险差 $P(y\mid x_1)-P(y\mid x_2)$——则只要 $Z$ 与 $X$ 无关联，或在给定 $X$ 时与 $Y$ 无关联，可压缩性便成立。因此，任何对可压缩性的违反都蕴含对定义 6.2.2 两个统计准则的违反；这大概就是许多人相信不可压缩性与混杂密切相关的缘故。然而，本章的例子表明，违反这两个条件对混杂而言既非充分也非必要。因此，不可压缩性与混杂一般是两个不同的概念；二者互不蕴含。

有些作者倾向于认为，这一区别是非线性效应度量 $g$（如优势比或似然比）的奇特性质，并且“当效应度量是总体单元上的期望时，混杂与不可压缩性在代数上等价”（Greenland 1998, p.~906）。本章表明，即便在线性泛函中，混杂与不可压缩性也未必对应。例如，效应度量 $P(y\mid x_1)-P(y\mid x_2)$（风险差）在图 6.3 中关于 $Z$ 不可压缩（对图的几乎每个参数化皆然），而该效应度量却是无混杂的（对每个参数化皆然）。

混杂与可压缩性之间的逻辑联系，是通过稳定无混杂的概念形成的，如定义 6.4.2 与定理 6.4.4 所表述。由于任何对可压缩性的违反都意味着违反定义 6.2.2 的 (U\textsubscript{1}) 与 (U\textsubscript{2})，它也就（由定理 6.4.4）蕴含违反稳定无偏性（或稳定无混杂）。于是我们可以陈述下述推论。

\begin{corollary}[稳定无混杂蕴含可压缩性]
设 $Z$ 是任何不受 $X$ 影响、且可能影响 $Y$ 的变量。设 $g[P(x,y)]$ 是任何度量 $X$ 与 $Y$ 之间关联的线性泛函。若 $g$ 关于 $Z$ 不可压缩，则 $X$ 与 $Y$ 不是稳定无混杂的。
\end{corollary}

这一推论为如下流行做法提供了理由：用参数变化法检验混杂，即只要“粗”关联度量 $g[P(x,y)]$ 不等于按 $Z$ 各水平平均的分 $Z$ 关联度量，就把变量 $Z$ 标为混杂因子（Breslow and Day 1980; Kleinbaum et al. 1982; Yanagawa 1984; Grayson 1987）。定理 6.4.4 提示，促成这一做法的直觉，源自对无混杂之稳定条件（而非仅仅偶然条件）的追求。此外，定理 6.4.4 中的条件 $A_Z$ 为某些作者提出的要求提供了理由：混杂因子必须是结局变量 $Y$ 的因果决定因素，而不仅仅与它关联。

\subsection{混杂与混杂因子}

本章讨论的焦点是混杂现象，我们把它等同于效应偏倚（定义 6.2.1）。关于这一主题的许多文献关注的是混杂因子的有无，预设某些变量具有混杂的能力而另一些没有。若照字面理解，这一观念可能引起误解；在把一个变量标为混杂因子之前，应当谨慎行事。

例如，Rothman and Greenland (1998, p.~120) 给出这样的定义：“造成暴露者与非暴露者疾病频率差异的额外因素，称为混杂因子”；他们接着指出：“一般地，一个混杂因子必须同时与研究中的暴露和研究中的疾病关联，才成其为混杂”（p.~121）。Rothman and Greenland 用“一般地”来限定其陈述，理由充分：我们已（在 \S 6.3.1 的双硬币例子中）看到，问题中的每个单独变量都可以同时与研究中的暴露（$X$）和疾病（$Y$）无关联，而 $X$ 对 $Y$ 的效应仍然是混杂的。类似情形也见于图 6.5 所描绘的线性模型。尽管 $Z$ 显然是 $X$ 对 $Y$ 之效应的混杂因子，因而必须加以控制，但 $Z$ 与 $Y$ 之间的关联实际上可能消失（在 $X$ 的每个水平上）；同样，$Z$ 与 $X$ 之间的关联也可能消失。若经由路径 $Z\leftarrow A\to Y$ 中介的间接关联恰好抵消了由箭头 $Z\to Y$ 承载的直接关联，这种情况就会发生。这种抵消并不意味着混杂不存在，因为路径 $X\leftarrow E\to Z\to Y$ 未被阻断，而 $X\leftarrow E\to Z\leftarrow A\to Y$ 被阻断。因此，$Z$ 是一个与疾病（$Y$）无关联的混杂因子。

\begin{figure}[htbp]
\centering
\begin{tikzpicture}[causal,scale=0.95,every node/.style={transform shape}]
\node[vtx] (E) at (1.2,2.6) {$E$};
\node[vtx] (A) at (5.0,2.6) {$A$};
\node[vtx] (Z) at (3.1,1.5) {$Z$};
\node[vtx] (X) at (1.0,0) {$X$};
\node[vtx] (Y) at (5.2,0) {$Y$};
\draw[ed] (E)--(X);
\draw[ed] (E)--(Z);
\draw[ed] (A)--(Z);
\draw[ed] (A)--(Y);
\draw[ed] (X)--(Y);
\draw[ed] (Z)--(Y);
\end{tikzpicture}
\caption{$Z$ 可能与 $X$ 和 $Y$ 都无关联，却仍是混杂因子（即每个充分集的成员）。}
\label{fig:6.5}
\end{figure}

Rothman and Greenland 刚刚所引陈述背后的直觉，可以通过稳定性概念形式地阐明：一个与 $X$ 或 $Y$ 稳定无关联的变量，可以安全地排除在调整之外。或者，Rothman and Greenland 的陈述也可以（不援引稳定性而）用非平凡充分集的概念（\S 3.3）来支持——所谓充分集，即对其作调整将消除混杂偏倚的变量集。可以证明（见本节末），每个这样的集合 $S$，作为一个整体，确实必须与 $X$ 关联，并在给定 $X$ 时与 $Y$ 条件关联。因此，Rothman and Greenland 的条件对非平凡充分（即可容许）集成立，但对集合中的个体变量不成立。

这一条件的实际含义如下。若给定一个声称充分（足以通过调整消除偏倚）的变量集 $S$，则该断言可以接受一个必要的统计检验：$S$ 作为一个复合变量，必须同时与 $X$ 关联、并与 $Y$（在给定 $X$ 时）关联。例如在图 6.5 中，$S_1=\{A,Z\}$ 与 $S_2=\{E,Z\}$ 都是充分且非平凡的；二者都必须满足所述条件。

然而请注意，尽管这一检验可以筛掉某些声称充分但显然糟糕的集合 $S$，它与充分性或混杂无关；它检验的只是非平凡性，即对 $S$ 作调整会改变 $X$ 与 $Y$ 之间的关联。当我们找到一个非平凡集 $S$ 时，我们仍不能确定该关联起初就是无偏的（如图 6.3），也不能确定它在调整之后变得无偏。

\subsubsection*{必要性的证明}

为证明只要 $Z$ 代表非平凡充分集 $S$，(U\textsubscript{1}) 与 (U\textsubscript{2}) 就必被违反，考虑 $X$ 对 $Y$ 没有影响的情形。此时，混杂相当于 $X$ 与 $Y$ 之间存在不消失的关联。条件独立的一个著名性质称为\emph{收缩}（contraction，\S 1.1.5），它断言：(U\textsubscript{1})，即 $X\indep S$，连同充分性 $X\indep Y\mid S$，蕴含对非平凡性的违反，即 $X\indep Y$：
\[
X\indep S \quad\&\quad X\indep Y\mid S
\qquad\Longrightarrow\qquad X\indep Y.
\]
同样，条件独立的另一性质称为\emph{交}（intersection），它断言：(U\textsubscript{2})，即 $S\indep Y\mid X$，连同充分性 $X\indep Y\mid S$，也蕴含对非平凡性的违反，即 $X\indep Y$：
\[
S\indep Y\mid X \quad\&\quad X\indep Y\mid S
\qquad\Longrightarrow\qquad X\indep Y.
\]
因此，任何非平凡充分集 $S$ 都必定同时违反 (U\textsubscript{1}) 与 (U\textsubscript{2})。

然而请注意，交性质只对严格正的概率分布成立；这意味着若问题中某些变量之间存在确定性关系，Rothman–Greenland 条件可能被违反。这可以从一个简单的例子看出：$X$ 与 $Y$ 都与第三个变量 $Z$ 处于一一对应的函数关系中。显然，$Z$ 是一个非平凡充分集，然而它在给定 $X$ 时与 $Y$ 无关联；一旦我们知道 $X$ 的值，$Y$ 的概率便已确定，不会再因得知 $Z$ 的值而改变。

\subsection{可交换性与混杂的结构分析}

流行病学的学生痛切抱怨文献中处理混杂这一基本概念的方式令人困惑。少数作者承认这种混乱（如 Greenland and Robins 1986; Wickramaratne and Holford 1987; Weinberg 1993），并提出了可能导向更系统分析的新视角。尤其是 Greenland 与 Robins（GR），他们认识到了我们在此 \S 6.2 与 \S 6.3 中所阐述的同样一些基本原理与结果。他们的分析是浩瀚混杂文献中少数亮点之一，因为它把混杂视为一种无法从观测数据直接测量的未知因果量。他们还承认（Miettinen and Cook 1981 亦然）：混杂的有无不应等同于可压缩性的无有，混杂也不应被视为依赖于参数的现象。

然而，本章所呈现的结构分析与 GR 的分析有一种根本区别：GR 走的是基于“可交换性”判断的进路。在 \S 6.1 中我们遇到过一种相关的可交换性概念，Lindley and Novick (1981) 曾试图用它来看待辛普森悖论；GR 的可交换性观念更具体、也更清晰地可应用。在概念上，混杂与可交换性之间的联系如下。若我们着手评估某种处理的效应，就应当确保处理组与未处理组之间的任何响应差异都归因于处理本身，而不是归因于两组之间与处理无关的某些内在差异。换言之，两组必须在一切与响应变量有关的特征上彼此相似。原则上，我们本可以在此结束混杂的定义，直接宣布：若处理组与未处理组在一切相关特征上彼此相似，则处理的效应无混杂。然而这一定义过于言辞化，因为它对“相似”与“相关”这两个词的解释高度敏感。为使其不那么非形式，GR 借用了 De Finetti 的假设性置换手法；研究者不再判断两组是否相似，而被要求设想两组的一次假设性交换（处理组变为未处理，反之亦然），然后判断交换之下的观测数据是否可与实际数据相区分。

人们有理由问：相对于直接判断两组是否实际上等同，这一心智练习究竟带来了什么？收获是双重的。第一，人很善于设想动态过程，能够根据对支配处理响应之过程、以及影响处理选择之因素的基本理解，模拟这一交换场景的结果。第二，从关于相似性的判断转向关于概率的判断，使我们能把这些判断写成概率记号，从而借用概率演算的力量与体面。

Greenland 与 Robins 迈出了通向这种形式化的重要第一步：把记号拉近判断的源头——人类对因果过程的理解。本书所走的结构进路则迈出下一个自然的步骤：把因果过程本身形式化。

设 $A$ 与 $B$ 分别代表处理组与未处理组，$P_{A1}(y)$ 与 $P_{A0}(y)$ 分别代表 $A$ 组在两种假设条件（处理与不处理）下的响应分布。\footnote{用 $\dox{\cdot}$ 记号，我们会写作 $P_{A1}(y)=P_A(y\mid \dox{X=1})$。}若我们关心响应分布的某个参数 $\mu$，则以 $\mu_{A1}$ 与 $\mu_{A0}$ 记该参数在相应分布 $P_{A1}(y)$ 与 $P_{A0}(y)$ 中的取值；对 $B$ 组，$\mu_{B1}$ 与 $\mu_{B0}$ 类似地定义。实际上，我们测量的是数对 $(\mu_{A1},\mu_{B0})$；在假设性交换之后，我们将测量 $(\mu_{B1},\mu_{A0})$。若这两个数对不可区分，即若
\[
(\mu_{A1},\mu_{B0}) = (\mu_{B1},\mu_{A0}),
\]
则称两组关于参数 $\mu$ 是\emph{可交换的}。

特别地，若我们用差值 $\mathrm{CE}=\mu_{A1}-\mu_{A0}$ 定义因果效应，则可交换性允许我们把 $\mu_{A0}$ 替换为 $\mu_{B0}$，从而得到 $\mathrm{CE}=\mu_{A1}-\mu_{B0}$；这是可测的，因为两个量都被观测到。Greenland 与 Robins 于是宣布：若 $\mu_{A0}=\mu_{B0}$，则因果效应 CE 无混杂。

若把这一定义与 \textup{(6.10)} 的定义 $P(y\mid \dox{x})=P(y\mid x)$ 相比较，我们会发现：若把后者改写为
\[
\mu[P(y\mid \dox{x})] = \mu[P(y\mid x)],
\]
其中 $\mu$ 是响应分布中所关心的参数，二者便重合。然而，结构进路与 GR 进路的主要差别在于分析的层次。结构建模在两个重要方向上扩展了混杂的形式化。第一，\textup{(6.10)} 不交由人的直接判断，而是从关于因果过程的更基本判断中数学地导出。\footnote{回想 $\dox{\cdot}$ 算子在结构方程模型中是以方程删除的方式数学地定义的；因此，在给定模型中验证无混杂条件 $P(y\mid \dox{x})=P(y\mid x)$，不是判断问题，而是数学分析的对象。}第二，结构模型所需的输入判断既是定性的，又是稳定的。

一个简单的例子将说明这些特征的益处。考虑下述陈述（Greenland 1998）：

\begin{quote}
(Q*)\quad “若效应度量是响应比例之差或之比，则上述现象——不可压缩而无混杂——不可能发生；混杂也不可能在不可压缩性缺席时发生。”（pp.~905--6）
\end{quote}

我们在本章中已经看到，陈述 (Q*) 应在若干方面加以限定；并且一般地，不可压缩性与混杂是两个不同的概念——无论采用何种效应度量，二者互不蕴含（\S 6.5.1）。然而，我们此处想讨论的问题是方法论上的：何种形式体系适于验证、反驳或限定这类陈述？显然，由于 (Q*) 对所有情形作出了一般断言，一个反例就足以驳倒其一般有效性。但我们如何构造这样的反例？更一般地，我们如何构造体现混杂、效应偏倚、因果效应、实验数据与非实验数据之别、反事实以及其他基于因果性之概念的例子？

在概率论中，若我们想反驳一个关于参数及其关系的一般陈述，只需给出一个使该关系不成立的密度函数 $f$。在命题逻辑中，为表明一个句子为假，只需给出一个满足前提却违反结论的真值表 $T$。那么，当我们想反驳像 (Q*) 这样的因果断言时，应当用什么数学对象来取代 $f$ 或 $T$？在 Greenland 与 Robins 的可交换性框架中，所用的相应对象是反事实列联表（例如见 Greenland et al. 1999b, p.~905，或 \S 1.4.4 的图 1.7）。例如，为说明混杂，我们需要两张这样的表：一张描述处理组 $A$ 对处理与不处理二者的假设性响应，另一张描述未处理组 $B$ 对处理与不处理二者的假设性响应。若表中显示：由处理组对不处理之假设性响应算出的参数 $\mu_{A0}$，不同于由未处理组之实际响应算出的 $\mu_{B0}$，则我们手上就有混杂。

这类表可以为涉及一个处理与一个响应变量的简单问题构造，但一旦涉及若干协变量，或当我们想对那些协变量施加某些约束时，它们就变成一场噩梦。例如，我们可能想纳入如下标准假设：协变量 $Z$ 不位于处理与响应之间的因果通路上，或者 $Z$ 对 $Y$ 有因果影响；但这类假设无法方便地在反事实列联表中表达。结果，被反驳之断言的原作者总能争辩说：反例中所用的表可能与公认的假设不相容。\footnote{尝试用反事实列联表为陈述 (Q*) 构造反例的读者，必定会体会到这一困难。}

这类困难不会困扰混杂的结构表示。在这一形式体系中，用以例示或反驳因果陈述的恰当对象是因果模型，如第三章所定义并在全书中所使用的那样。在这里，假设性响应 $\mu_{A0}$ 与 $\mu_{B0}$、以及列联表，都不是原初量，而是可从一组已经体现我们所愿遵守之假设的方程中导出的。结构模型的每个参数化（用 \textup{(3.51)} 或 $\dox{\cdot}$ 算子）都蕴含一组特定的反事实列联表，它满足输入假设，并展示图中所示的统计性质。例如，图 6.3 之图的任何参数化都生成一组反事实列联表，其中已经体现“$Z$ 不在 $X$ 与 $Y$ 之间的因果通路上”以及“$Z$ 对 $Y$ 没有因果效应”这些假设；而几乎每个这样的参数化都会生成 (Q*) 的一个反例。此外，我们也可以仅凭对图的随意一瞥、而无须生成数值反例，就否证 (Q*)。例如，图 6.3 生动地表明，风险差 $P(y\mid x_1)-P(y\mid x_2)$ 关于 $Z$ 不可压缩，同时 $X$ 与 $Y$ 是（稳定）无混杂的。

当我们不是要反驳、而是要证实关于混杂的一般断言时，两种表述之间的差别更为显著。此时，仅呈现一张列联表是不够的；我们必须证明该断言对一切可能按输入假设构造出来的表都成立。读者必定明白，在列联表框架内这是一项无望的工作；它要求一种形式体系，其中假设可以被简洁地陈述，结论可以通过数学推导得出。结构语义提供了这样的形式体系，正如本书中证明的许多一般断言所展示的那样（例子包括定理 6.4.4 与推论 6.5.2）。

尽管我十分钦佩 Greenland 与 Robins 通过可交换性框架引入的严格性，我仍深信：反事实列联表的不透明与僵硬，在很大程度上造成了 GR 框架在流行病学家中接受缓慢，并且作为副产品，造成了统计文献中围绕混杂的挥之不去的混乱。我同样深信，用结构模型的语言表述断言与假设，将使因果性的数学分析为普通研究者所掌握，从而最终导致对混杂的一次彻底而自然的“去混杂”。

\section{结论}

过去为在统计关联（或可压缩性）与混杂之间建立理论联系的努力之所以未获成功，原因有三。第一，缺乏表达关于因果关系与效应偏倚之断言的数学语言，使得人们难以评估效应无偏性要求（定义 6.2.1）与声称捕捉无偏性的统计准则之间的差距。\footnote{关于可压缩性的大多数论文（如 Bishop 1971; Whittemore 1978; Wermuth 1987; Becher 1992; Geng 1992）都通过援引辛普森悖论以及获得混杂效应估计的危险来引出主题。其中只有少数研究混杂或效应估计；多数宁愿把较易处理的不可压缩性现象作为独立目标来分析。有些甚至把可压缩性命名为“无混杂性”（Grayson 1987; Steyer et al. 1997）。}第二，把贫瘠代理（图 6.3）排除出考量的必要性，在某种程度上未引起研究者的注意。最后，稳定无偏与偶然无偏之间的区别未得到应有的重视；而且正如我们在例 6.3.3 中所见，若不承诺稳定性概念，就无法在关联准则（或可压缩性）与混杂之间建立任何联系。这种承诺关键地依赖于如下观念：因果模型是由若干可以彼此独立变化的自主机制装配而成的（Aldrich 1989）。正是出于对这种独立变化的预期，我们才不满足于偶然无偏，而寻求稳定无偏的条件。这一观念在数学上的形式化，导致了 DAG-同构（Pearl 1988b, p.~128）、稳定性（Pearl and Verma 1991）与忠实性（Spirtes et al. 1993）等相关概念，它们有助于从稀疏的统计关联中阐明因果图（见第二章）。那些渴望把关联准则与混杂联系起来的作者，显然也共享同一观念。

结构模型分析连同图形方法的辅助，提供了一个数学框架，使混杂的考量能够被更有效地表述与管理。利用这一框架，本章阐明了稳定无偏性准则，并表明该准则 (i) 一直是流行病学与生物统计学中许多研究暗中追求的目标，(ii) 可以给出类似于检验可压缩性时所用那样的操作统计检验。我们进一步表明（\S 6.5.3），结构框架克服了使混杂成为文献中最混乱主题之一的基本认知与方法论障碍。因此，可以自然地预言：这一框架将成为未来混杂研究的主要数学基础。

\section*{致谢}
\addcontentsline{toc}{section}{致谢}

与 James Robins 和 Sander Greenland 的讨论极有价值。尤其是 Sander，他对本书的两份早期草稿提出了许多建设性意见，并帮助使其对流行病学家保持可理解。Jan Koster 提请我注意 Stone 与 Robins 的无混杂准则之间的联系，并指出了早期草稿中的若干疏漏。其他有益的讨论者还有 Michelle Pearl、Bill Shipley、Rolf Steyer、Stephen Stigler 与 David Trichler。

\section*{第二版后记}
\addcontentsline{toc}{section}{第二版后记}

读者会好笑地得知，本章的第一次印刷并未终止统计学家对辛普森非悖论的迷恋。教科书继续对这一现象惊叹不已（Moore and McCabe 2005），研究者继续以 1970—90 年代的热情追逐其数学精妙（Cox and Wermuth 2003）与可视化（Rücker and Schumacher 2008），却一次也不提“原因”一词，也一次不停下来追问：“意义何在？”一个显著的例外是 Larry Wasserman 的《All of Statistics》（Wasserman 2004），它是第一本在正确的因果语境中处理辛普森逆转的统计学教科书。我对这一非悖论之因果理解终将获胜的信心，也因流行病学文献中发表的一场讨论而得到加强，该讨论以毫不含糊的措辞总结道：“解释与解决方案在于依赖背景知识、而非统计准则的因果推理”（Arah 2008）。看来，当我们进入因果性的时代，我们应当向流行病学家寻求指引与智慧。
